\chapter{Mathieu Moonshine Computed: Twining, the Shadow, and the Little-String Bridge}\label{ch:moonshinecomputed}

In \cref{ch:moonshine} the coincidence between the elliptic-genus multiplicities
$2A_n$ and the dimensions of $\Mtf$ was checked in Lean as one typed table against
another. That is weak evidence: a table matching a table only shows that two lists of
numbers were copied consistently, and this project has already caught one mistyped
$\Mtf$ table. This chapter asks a sharper question. \emph{Can both sides of Mathieu
moonshine be recomputed from formulas---the $q$-series from modular forms, the traces from
the character table---so that the agreement is a computation and not a transcription? And
which integers of this programme survive the test that separates moonshine from
numerology?}\index{Mathieu moonshine!computed} The answers are: yes, at every one of the
$26$ conjugacy classes and for the first ten graded pieces; the programme's own ``$27720$
lock'' fails the test at every non-trivial class, while the number $24=\chi(\mathrm{K3})$
passes it; and the same arithmetic reproduces a bridge, due to Harvey--Murthy--Nazaroglu,
between the elliptic genus and the BPS states of double-scaled little string theories.
Along the way the student learns a technique worth keeping: treating a truncated
$q$-series as an exact list of integers, so that a proof assistant can decide an
identity between modular objects coefficient by coefficient.

Conventions follow \cref{ch:moonshine}: $q=\ee^{2\pi\ii\tau}$, $y=\ee^{2\pi\ii z}$,
$A_n=45,231,770,\dots$ are the EOT coefficients \source{EOT eq.~(1.12), l.~208} and the
massive multiplicities are $2A_n$. The source keys CDH (Cheng--Duncan--Harvey,
\emph{Umbral Moonshine}), CH (Cheng--Harrison), HMN (Harvey--Murthy--Nazaroglu) and Cheng
are listed in the bibliography.

\section{Why compute instead of type}\label{sec:moonshinecomputed:why}

A moonshine statement has two sides. On the \emph{modular} side there is a function,
here the mock modular form
\begin{equation}\label{eq:moonshinecomputed:H}
  H(\tau)=2q^{-1/8}\bigl(-1+45q+231q^2+770q^3+\dots\bigr),
\end{equation}
whose coefficients are fixed by modularity and a few pieces of data. On the \emph{group}
side there is a graded representation $K=\bigoplus_nK_n$ of $\Mtf$ whose dimensions are
those coefficients. The literature records both sides as tables: the $A_n$ as printed
numbers, the $\Mtf$ representations through a printed character table. Checking one
table against the other tests the copying, not the mathematics.

The alternative is to compute each side from a short formula and compare the outputs.
Three things then become checkable that were not before. First, a misprint in a table can
no longer hide: the computed column either equals the printed one or it does not. Second,
one can compare not only dimensions but \emph{traces} of every group element, which is
the real content of moonshine. Third, one can apply the same test to any other integer
coincidence, including this programme's own. The price is that one needs an exact
arithmetic for power series; the next section supplies it.

\section{Truncated \texorpdfstring{$q$}{q}-series as lists of integers}\label{sec:moonshinecomputed:series}

A power series $f=\sum_{n\ge0}f_nq^n$ with integer coefficients, known through
$q^N$, is stored as the list $[f_0,f_1,\dots,f_N]$ of $N+1$ integers.\index{q-series@$q$-series!truncated}
Two operations suffice.\index{Cauchy product!truncated} The truncated Cauchy product
\begin{equation}\label{eq:moonshinecomputed:mul}
  (fg)_n=\sum_{k=0}^{n}f_kg_{n-k},\qquad 0\le n\le N,
\end{equation}
is exact: the coefficients of $fg$ through $q^N$ depend only on those of $f$ and $g$
through $q^N$. If $p_0=1$ the quotient $r=s/p$ is computed by the recursion
\begin{equation}\label{eq:moonshinecomputed:div}
  r_n=s_n-\sum_{k=0}^{n-1}r_k\,p_{n-k},
\end{equation}
obtained by reading off the coefficient of $q^n$ in $s=rp$; it too is exact, and it
involves no division by anything but $p_0=1$. With these two operations every eta
product, every Eisenstein series and every quotient appearing below becomes a finite list
of integers, and an identity between two such lists is decided by comparing them entry by
entry. Nothing in this is approximate; everything is also \emph{finite}: an identity
checked through $q^9$ says nothing about $q^{10}$.

\subsection{Worked example: the first coefficients of \texorpdfstring{$H$}{H}}\label{sec:moonshinecomputed:worked}

Cheng and Harrison write the mock modular form of the Niemeier root system $A_1^{24}$
(the case of Mathieu moonshine) in closed form \source{CH eqs.~(3.5)--(3.6), ll.~797--815}:
\begin{equation}\label{eq:moonshinecomputed:CH}
  H(\tau)=\frac{-2E_2(\tau)+48F_2(\tau)}{\eta(\tau)^3},\qquad
  F_2(\tau)=\sum_{\substack{r>s>0\\ r-s\ \mathrm{odd}}}(-1)^r s\,q^{rs/2}
  =q+q^2-q^3+q^4+\dots,
\end{equation}
with the Eisenstein series\index{Eisenstein series!$E_2$} $E_2(\tau)=1-24\sum_{n\ge1}\sigma_1(n)q^n$ \source{HMN eq.~(2.37), l.~736}. Since
$\eta^3=q^{1/8}\prod_{n\ge1}(1-q^n)^3$, the series $q^{1/8}H$ is the quotient of the
integer series
\begin{equation}
  s=-2E_2+48F_2=-2+48\bigl(\sigma_1(n)+[q^n]F_2\bigr)_{n\ge1}
  =-2+96q+192q^2+144q^3+\dots
\end{equation}
by $p=\prod(1-q^n)^3=1-3q+5q^3-7q^6+\dots$ (Jacobi's identity; in Lean it is checked
through $q^9$, not proved). The recursion \eqref{eq:moonshinecomputed:div} gives
\begin{align*}
  r_0&=-2,\\
  r_1&=96-r_0p_1=96-(-2)(-3)=90,\\
  r_2&=192-(r_0p_2+r_1p_1)=192-(0-270)=462,\\
  r_3&=144-(r_0p_3+r_1p_2+r_2p_1)=144-(-10+0-1386)=1540.
\end{align*}
These are $-2,\,2A_1,\,2A_2,\,2A_3$. Carried through $q^9$, the same recursion reproduces
all nine tabulated values $2\cdot45,\dots,2\cdot132825$; in Lean this is the theorem
\lean{hComputed\_eq\_table} of \cref{sec:moonshinecomputed:lean}. \TA\ for the equality
of lists; \TL\ for the statement that \eqref{eq:moonshinecomputed:CH} is the mock modular
form attached to the K3 elliptic genus.

\begin{pitfallbox}[Ten coefficients are not an identity]
The theorem \lean{eta3\_jacobi} says that $\prod(1-q^n)^3$ and Jacobi's series
$\sum(-1)^n(2n+1)q^{n(n+1)/2}$ agree through $q^9$. That is a check of ten numbers, not a
proof of Jacobi's triple product identity, and the chapter never uses it as one. The same
holds for every ``$=$'' in this chapter: it is an equality of truncations.
\end{pitfallbox}

\section{Twining from the modular side}\label{sec:moonshinecomputed:twining}

Moonshine becomes structural when dimensions are replaced by traces\index{twining genus}
(\cref{ch:moonshine}). For each conjugacy class $g$ of $\Mtf$ the twined series $H_g$
has coefficients $\Tr_{K_n}(g)$. Cheng--Duncan--Harvey give all of them in one formula
\source{CDH eq.~(4.18), ll.~2931--2936}:
\begin{equation}\label{eq:moonshinecomputed:twined}
  H_g(\tau)=\frac{\chi_g}{24}\,H(\tau)+\frac{F_g(\tau)}{\eta(\tau)^3},
\end{equation}
where $\chi_g$ is the number of points of the $24$-point set fixed by $g$
\source{CDH Table~14, ll.~4349--4364} and $F_g$ is a weight-$2$ modular form for
$\Gamma_0(N_g)$ listed in \source{CDH Table~3, ll.~2882--2915}. Most $F_g$ are built from
$\Lambda_N=Nq\,\partial_q\log\bigl(\eta(N\tau)/\eta(\tau)\bigr)$ \source{CDH eq.~(A.3),
l.~4100}, which is an integer series after multiplication by $24$:
\begin{equation}
  24\Lambda_N=N(N-1)+24N\sum_{n\ge1}\Bigl(\sigma_1(n)-N\sigma_1(n/N)\Bigr)q^n,
\end{equation}
the second term present only when $N\mid n$. To stay in $\Z$ one computes
$24H_g=(\chi_g\,s+24F_g)/p$ with $s$ and $p$ as in \cref{sec:moonshinecomputed:worked}.

\begin{example}[The class $2A$]
Here $\chi_{2A}=8$ and $F_{2A}=-16\Lambda_2$ \source{CDH Table~3, l.~2886}, so
$24F_{2A}=-16\cdot24\Lambda_2=-32-768q+\dots$. The numerator is
$8s+24F_{2A}=(-16-32)+(768-768)q+\dots=-48+0\cdot q+\dots$, and the recursion gives
$r_0=-48$, $r_1=0-(-48)(-3)=-144$. Dividing by $24$: $-2,\,-6$, the first two entries of
the $2A$ column of \source{CDH Table~20, ll.~4494--4512}. The trace of $2A$ on the
$90$-dimensional $K_1=\mathbf{45}\oplus\overline{\mathbf{45}}$ is therefore $-6$, which is
$2\chi_{\mathbf{45}}(2A)=2\cdot(-3)$ from the character table (\cref{ch:m24}).
\end{example}

\subsection{Eta quotients and their shift}\label{sec:moonshinecomputed:etaq}

Several $F_g$ are eta quotients,\index{eta quotient} and these carry a fractional power of $q$ that has to be
accounted for before anything is stored in a list of integers. Since
$\eta(k\tau)=q^{k/24}\prod_{n\ge1}(1-q^{kn})$, a quotient
$\prod_k\eta(k\tau)^{e_k}$ equals
\begin{equation}\label{eq:moonshinecomputed:etaq}
  q^{\,\sum_kke_k/24}\prod_k\prod_{n\ge1}(1-q^{kn})^{e_k},
\end{equation}
where the second factor is an integer series computed by
\eqref{eq:moonshinecomputed:mul}--\eqref{eq:moonshinecomputed:div}. The prefactor is a
genuine power of $q$ exactly when $24\mid\sum_kke_k$, and each quotient used here satisfies
that. For the newform $f_{11}=\eta(\tau)^2\eta(11\tau)^2$ one has
$\sum_kke_k=2\cdot1+2\cdot11=24$, so $f_{11}=q\prod_n(1-q^n)^2(1-q^{11n})^2$, and indeed
$f_{11}=q-2q^2-q^3+2q^4+\dots$ begins at $q^1$, as the normalisation
$f_{23,a}=q+O(q^3)$, $f_{23,b}=q^2+O(q^3)$ of \source{CDH eq.~(A.8), ll.~4134--4137}
requires for its companions. The Lean file checks the divisibility condition for every
quotient it uses before using it, and checks these normalisations; without them a shift of
one power of $q$ would silently move a whole column.

\begin{pitfallbox}[A displaced exponent]
In the text extraction of CDH's Table~3 the exponent $3$ of $\eta(\tau)$ in the $12A$ entry
sits on the next line, so the entry reads as if the exponent were missing. It was
reconstructed by requiring weight $2$ (the exponents of an eta quotient of weight $w$ sum to
$2w$), and the reconstruction is confirmed by the match with the printed column; a Lean
control shows that the exponent $2$ instead of $3$ does \emph{not} reproduce it. This is the
kind of error a computed check catches and a typed table hides.
\end{pitfallbox}

Sixteen of the $21$ columns of CDH's Table~20 need more than a single $\Lambda_N$: eta
quotients, sums of several $\Lambda_d$, and for $11A$, $14AB$, $15AB$ and $23AB$ the
newforms of CDH's Appendix~A with rational prefactors. For instance
$F_{11A}=\tfrac25(-\Lambda_{11}+11f_{11})$ \source{CDH Table~3, l.~2900} with the newform
$f_{11}=\eta(\tau)^2\eta(11\tau)^2$ \source{CDH eq.~(A.4), l.~4115}. An eta quotient
$\prod_k\eta(k\tau)^{e_k}$ is $q^{\sum_kke_k/24}$ times an integer series, and every
quotient used here has $\sum_kke_k\equiv0\pmod{24}$, so the power of $q$ is an integer.
A prefactor $1/D$ is handled by computing $24D\,H_g$. In Lean all $21$ computed columns
equal the printed ones through $q^9$ (\TA), and each has every coefficient divisible by
$24D$, so the division back to $H_g$ is exact. Two negative controls make the agreement
informative: without $F_g$ the $2A$ column is not reproduced, and without the newform
$f_{11}$ the $11A$ column is not either. Because CDH's table was read from a flattened
PDF, these theorems also certify the transcription; one exponent in CDH's Table~3 is
displaced in the text extraction ($12A$) and was fixed by weight counting, which the
match then confirms.

\section{Twining from the group side}\label{sec:moonshinecomputed:group}

\subsection{A character table with irrational entries}

The other side needs the traces $\Tr_{K_n}(g)$ computed from the character table of
$\Mtf$ \source{CDH Table~8, ll.~4170--4209} and the decompositions of the $K_n$. Most
entries are integers; a few are the algebraic integers\index{character table!of M24@of $\Mtf$}
\begin{equation}
  b_n=\frac{-1+\sqrt{-n}}{2},\qquad n=7,15,23,\qquad b_n^2=-b_n-\frac{n+1}{4},\qquad
  \overline{b_n}=-1-b_n.
\end{equation}
An entry $a+b\,b_n$ is stored as the integer pair $(a,b)$; multiplication and complex
conjugation are the two polynomial rules just written, so all arithmetic stays in $\Z$.
Each column of the table lives in one ring: $\Z$, $\Z[b_7]$ (classes $7A,7B,14A,14B,21A,
21B$), $\Z[b_{15}]$ ($15A,15B$) or $\Z[b_{23}]$ ($23A,23B$).

\subsection{Guards: what a character table must satisfy}

A mistyped table is how this project went wrong before, so the transcription is guarded
by identities every character table satisfies. The column norms
$\sum_i|\chi_i(g)|^2$ equal the centralizer orders $|C(g)|$ (second orthogonality);
Lean computes them from the table, rather than taking them from memory, and checks that
each divides $|\Mtf|=244\,823\,040$ and is divisible by the order of $g$. The class sizes
$|\Mtf|/|C(g)|$ then add up to $|\Mtf|$ (the class equation).\index{class equation} Finally the first
orthogonality relation\index{orthogonality relations}
\begin{equation}\label{eq:moonshinecomputed:orth}
  \sum_g|g^{\Mtf}|\,\chi_i(g)\,\overline{\chi_j(g)}=|\Mtf|\,\delta_{ij}
\end{equation}
is checked for all $26\times26$ pairs. The sum mixes values in four rings; it is
accumulated as a quadruple (rational part, coefficient of $b_7$, of $b_{15}$, of
$b_{23}$). Since $1,\sqrt{-7},\sqrt{-15},\sqrt{-23}$ are linearly independent over $\Q$
(standard; \TL), \eqref{eq:moonshinecomputed:orth} holds exactly when the quadruple is
$(|\Mtf|\delta_{ij},0,0,0)$.

\begin{pitfallbox}[What the guards do not fix]
The text extraction of CDH's table loses the overlines: $b_7$ and $\overline{b_7}$ both
print as ``b7''. A search over the ten independent choices (in Python, not in the kernel)
finds $128$ readings that pass all guards---exactly the orbit of one reading under
renaming the classes $7A\leftrightarrow7B$, $14A\leftrightarrow14B$, \dots\ and the
characters $\mathbf{45}\leftrightarrow\overline{\mathbf{45}}$, \dots---and all of them give
the same traces below. So the guards pin the table up to relabelling, not absolutely, and
the chapter's statements are invariant under that relabelling. A reading \emph{outside}
the orbit is rejected: the Lean control \lean{misread\_fails\_gram} swaps one pair of
entries and orthogonality fails.
\end{pitfallbox}

\subsection{Traces equal computed coefficients}

With the table in hand, the trace of $g$ on $K_n=\bigoplus_im_{n,i}\rho_i$ is
$\sum_im_{n,i}\chi_i(g)$. For levels $n=1,\dots,7$ EOT give decompositions
(\cref{ch:moonshine}), and Lean checks that at \emph{every} class the resulting trace is a
rational integer (its $b_n$-parts cancel) and equals the coefficient of $q^{n-1/8}$ in the
computed $H_g$ (\TA). Two independent computations---one from modular forms, one from
the character table---meet at $26\times7$ numbers.

\begin{example}[The class $3A$ at levels $1,2,3$]
The $3A$ column of the character table gives $\chi_{\mathbf{45}}(3A)=0$,
$\chi_{\mathbf{231}}(3A)=-3$, $\chi_{\mathbf{770}}(3A)=5$ \source{CDH Table~8,
ll.~4170--4209}. Since $K_1=\mathbf{45}\oplus\overline{\mathbf{45}}$,
$K_2=\mathbf{231}\oplus\overline{\mathbf{231}}$ and
$K_3=\mathbf{770}\oplus\overline{\mathbf{770}}$, and each pair contributes twice the real
part, the traces are
\[
  \Tr_{K_1}(3A)=0,\qquad \Tr_{K_2}(3A)=-6,\qquad \Tr_{K_3}(3A)=10 .
\]
On the modular side, $\chi_{3A}=6$ and $F_{3A}=-6\Lambda_3$ \source{CDH Table~3,
ll.~2882--2915} give the column $-2,\,0,\,-6,\,10,\,0,\,-18,\dots$
\source{CDH Table~20, ll.~4494--4512}. The three numbers agree---and they agree for all
$26$ classes and all seven levels at once, which is the content of
\lean{trace\_eq\_twined\_coeff\_all}. Note how little room there is for an accident: the
modular side knows only $\chi_{3A}$ and one weight-two form, the group side only a column
of characters.
\end{example}

\begin{figure}[htbp]
\centering
\begin{tikzpicture}[>=stealth,node distance=6mm,
  box/.style={draw,rounded corners,align=center,inner sep=2.6mm,font=\small}]
  \node[box] (mod) at (-4.4,0)
    {\textbf{modular side}\\[1mm] $E_2,\ F_2,\ \Lambda_N,\ \eta$\\ newforms $f_{11},f_{14},f_{15},f_{23}$\\[1mm]
     $H_g=\frac{\chi_g}{24}H+\frac{F_g}{\eta^3}$};
  \node[box] (grp) at (4.4,0)
    {\textbf{group side}\\[1mm] character table over\\ $\Z,\ \Z[b_7],\ \Z[b_{15}],\ \Z[b_{23}]$\\[1mm]
     $\Tr_{K_n}(g)=\sum_i m_{n,i}\chi_i(g)$};
  \node[box,fill=black!4] (mid) at (0,0) {lists of integers\\ through $q^9$};
  \draw[->] (mod) -- node[above,font=\scriptsize]{computed} (mid);
  \draw[->] (grp) -- node[above,font=\scriptsize]{computed} (mid);
  \node[box,fill=tierA!12] (out) at (0,-2.1)
    {equal at all $26$ classes, levels $0$--$9$\\[0.5mm]
     \scriptsize (and: $27720$ ratio fails at $25$ classes; $\chi_g$ is a character)};
  \draw[->] (mid) -- (out);
\end{tikzpicture}
\caption{The two routes of this chapter. Neither side is copied from the other: the left
side is computed from modular forms, the right side from a guarded character table, and the
comparison is made on finite lists of integers.}
\label{fig:moonshinecomputed:tworoutes}
\end{figure}

\begin{physicsbox}[What the twined numbers count]
$\Tr_{K_n}(g)$ is a graded character: it counts the BPS states of the K3 sigma model at
level $n$ with a sign given by how the symmetry $g$ acts on them. Two consequences make it
a sharper probe than the dimension. Because characters are additive over direct sums, any
relation coming from a module structure survives the replacement of dimensions by traces.
And because a non-trivial $g$ mixes states, the traces are small numbers of both signs, so
an accidental numerical relation among the $\dim K_n$ has no reason to persist. This is why
the comparison in \cref{fig:moonshinecomputed:tworoutes} is evidence for a module and the
$27720$ relation of \cref{sec:moonshinecomputed:forger} is not.
\end{physicsbox}

For $n=8,9$ EOT give no decomposition. The direction can then be reversed: the
multiplicity of $\rho_i$ in $K_n$ is the inner product
\begin{equation}\label{eq:moonshinecomputed:mult}
  m_{n,i}=\frac1{|\Mtf|}\sum_g|g^{\Mtf}|\;[q^{n-1/8}]H_g\;\overline{\chi_i(g)}.
\end{equation}
Lean evaluates \eqref{eq:moonshinecomputed:mult} for $n=0,\dots,9$ and all $26$ irreducibles
from the \emph{computed} $H_g$: every $m_{n,i}$ is an integer, non-negative for $n\ge1$,
and equal to CDH's printed decomposition \source{CDH Table~48, ll.~5305--5330}.

\begin{example}[Level $8$]
The computed multiplicities give
\begin{align*}
  K_8=\ &990\oplus\overline{990}\oplus1035'\oplus1035''\oplus2\cdot1265\oplus2\cdot2277\\
  &\oplus\,2\cdot3312\oplus2\cdot3520\oplus4\cdot5313\oplus2\cdot5796\oplus2\cdot5544
   \oplus6\cdot10395,
\end{align*}
whose dimension is $1980+2070+2530+4554+6624+7040+21252+11592+11088+62370=131100
=2A_8$. The level-$7$ decomposition $10395\oplus5796\oplus5544\oplus5313\oplus2024\oplus1771$
(doubled) is the one used in \cref{ch:moonshine}; the earlier proposal
$10395+2\cdot5796+5544+3312$ has the right dimension but the wrong traces, as Cheng
observed \source{Cheng, ll.~533--549}, and Lean confirms it disagrees with the computed
twined series at some class.
\end{example}

\section{The forger's test}\label{sec:moonshinecomputed:forger}

Anyone can find integers that match: $A_2\cdot60=4A_1\cdot77=27720$, with
$27720\cdot8832=|\Mtf|$, is the ``lock'' of \cref{ch:moonshine}. What made monstrous
moonshine more than numerology was that the coincidences survive \emph{twining}: replace
every dimension by the trace of an element and the relation still holds.\index{forger's test}\index{twining!test}
Apply this to the lock in its ratio form,
$[q^2]H_g\cdot(4\cdot[q^1]H)=[q^2]H\cdot(4\cdot[q^1]H_g)$, i.e.\ $A_2(g)/A_1(g)=A_2/A_1$.
At $g=2A$ the computed coefficients are $A_1(2A)=-6$, $A_2(2A)=14$, and $14/(-6)\neq
462/90$. Lean checks that the relation holds at the identity and fails at every one of the
$25$ non-identity classes (\TA). By this criterion the lock is numerology.

The number $24$ behaves differently. The multiplicity $\chi_g$ in
\eqref{eq:moonshinecomputed:twined}, which replaces $24$ when one twines, is at every class
the trace of $g$ on the permutation representation $\mathbf1\oplus\mathbf{23}$ computed
from the character table: $24,8,0,6,0,0,4,\dots$ (Lean checks all $26$ entries). The
$24$ twines into a genuine character.

\begin{physicsbox}[Why twining separates structure from coincidence]
A relation that comes from a module structure---``$K_n$ is this sum of representations''---
holds for traces because traces are additive over direct sums. A relation between
dimensions that is only numerical has no reason to survive: the traces of a non-trivial
element are smaller, of both signs, and unrelated by any multiplicative identity. The test
is cheap once both sides are computed, and it should be applied to every integer
coincidence before a physical reading is proposed.
\end{physicsbox}

\section{Where the 24 of the shadow comes from}\label{sec:moonshinecomputed:shadow}

$H$ is a mock modular form: it becomes modular only after adding a non-holomorphic
correction built from its \emph{shadow}, which for Mathieu moonshine is
$24\eta(\tau)^3$ \source{HMN §1, ll.~244--248}.\index{shadow}\index{mock modular form!shadow}
The $24$ in the shadow can be located arithmetically. CDH write any weight-$1$
meromorphic Jacobi form of the shape $\psi=\Psi_{1,1}\phi$, with
$\Psi_{1,1}=-\ii\,\theta_1(\tau,2z)\eta^3/\theta_1(\tau,z)^2$ and $\phi$ a weak Jacobi form
of weight $0$, as a polar part plus a finite part \source{CDH eqs.~(2.19)--(2.25),
ll.~657--711}:
\begin{equation}\label{eq:moonshinecomputed:polar}
  \psi=\chi\,\mathrm{Av}^{(m)}\!\Bigl[\frac{y+1}{y-1}\Bigr]+\sum_rh_r\,\hat\theta_r^{(m)},
  \qquad \chi=\phi(\tau,0),
\end{equation}
and the shadow of $(h_r)$ is $\chi$ times the unary theta series\index{theta series!unary} $S^{(m)}$
\source{CDH eqs.~(2.26)--(2.30), ll.~716--768}. For K3, $\phi$ is the elliptic genus
$Z=8(f_2^2+f_3^2+f_4^2)$ with $f_i=\theta_i(\tau,z)/\theta_i(\tau,0)$ \source{CDH
eqs.~(2.49), (2.62), ll.~1258--1270, 1397}, $m=2$, and $S^{(2)}_1=\eta^3$.

Everything in \eqref{eq:moonshinecomputed:polar} except the word ``shadow'' is a finite
computation once denominators are cleared: the theta functions are products, $Z$ is a
$q$-series whose coefficients are Laurent polynomials in $y$, and the Appell--Lerch sum
$\mathrm{Av}^{(2)}$ is a sum of geometric series. Lean checks three facts through $q^9$.
First, $Z(\tau,0)=24$: the constant term of $Z$ at $y=1$ is $24$ and every higher
coefficient vanishes. Second, the decomposition \eqref{eq:moonshinecomputed:polar} holds
with $\chi=24$ and $h_1$ equal to the series $H$ \emph{computed} from
\eqref{eq:moonshinecomputed:CH}. Third, it fails with $\chi=23$ or $\chi=25$. So two
independent facts meet at $24$: the value $Z(\tau,0)$, i.e.\ the Euler characteristic
of K3 (the Witten index), and the multiplicity of the polar part. That this multiplicity is
the coefficient of the shadow is Zwegers' theorem as recast by CDH (\TL); nothing in Lean
constructs the non-holomorphic completion.

\begin{pitfallbox}[Located, not proved]
The Lean files never define a shadow, a completion or a modular transformation. What is
proved is the arithmetic skeleton: the $24$ of $Z(\tau,0)$ and the $24$ of the polar part
are the same integer, and no other integer makes the decomposition hold. ``$H$ is mock
modular with shadow $24\eta^3$'' remains \TL.
\end{pitfallbox}

\section{The little-string bridge}\label{sec:moonshinecomputed:hmn}

\subsection{A BPS index labelled by ADE diagrams}

Harvey, Murthy and Nazaroglu study $k$ NS5-branes on K3 in a double-scaling limit; the
branes' positions define an ADE root system $Y$ with Coxeter number $k$, the lattice of
vanishing $(-2)$-cycles.\index{little string theory!double-scaled}\index{ADE} The holomorphic part of the second
helicity supertrace\index{helicity supertrace}, a BPS index of this theory, is given in closed form
\source{HMN eqs.~(2.29)--(2.37), ll.~700--736}:
\begin{equation}\label{eq:moonshinecomputed:chi2Y}
  \chi_2^Y=\mathrm{rk}(Y)\,E_2-24\,F_2^Y,\qquad
  F_2^{(k,d)}=\sum_{n\ge1}q^n\sum_{rs=n}s\bigl(d\,[kr>d^2s]-\tfrac kd\,[d^2r>ks]\bigr),
\end{equation}
with $F_2^Y$ a combination of $F_2^{(k,d)}$ fixed by $Y$ ($A_{k-1}$: $F_2^{(k,1)}$;
$D_{m/2+1}$: $F_2^{(m,1)}+F_2^{(m,m/2)}$; and three combinations for $E_{6,7,8}$).

\begin{example}[$Y=A_3$]
Here $k=4$, $\mathrm{rk}=3$, and at $n=1$ the only factorization is $r=s=1$, giving
$F_2^{(4,1)}(1)=1\cdot(1\cdot[4>1]-4\cdot[1>4])=1$. Hence
$[q^1]\chi_2^{A_3}=3\cdot(-24)-24\cdot1=-96$, the second coefficient of HMN's printed
series $3-96q-288q^2-384q^3-576q^4-360q^5+\dots$ \source{HMN eq.~(4.10), l.~1290}. Lean
reproduces all six printed coefficients, and the six of $Y=A_7$ \source{HMN eq.~(4.11),
l.~1292}, whose coefficients $7$ and $192$ share no factor.
\end{example}

For $k=2$ ($Y=A_1$, two NS5-branes) the formula gives $-2\chi_2^{A_1}=-2E_2+48F_2$, the
numerator of \eqref{eq:moonshinecomputed:CH}: the little-string index is
$\eta^3H$, the weight-$2$ form of Mathieu moonshine (Lean checks it through $q^9$, and that
HMN's sum $F_2^{(2,1)}$ equals Cheng--Harrison's $F_2$ through $q^{40}$).

\subsection{The umbral relation and a divisibility}

\begin{example}[The relation \eqref{eq:moonshinecomputed:umbral} at $\ell=2$, by hand]
Take $Y=A_1$, so $k=2$, $\mathrm{rk}(Y)=1$ and $X=A_1^{24}$, whose umbral form is the $H$ of
\eqref{eq:moonshinecomputed:CH}; here $\chi_2^X=\eta^3H=-2E_2+48F_2$, whose $q^1$
coefficient is $96$ (\cref{sec:moonshinecomputed:worked}). The right-hand side of
\eqref{eq:moonshinecomputed:umbral} is then $-\tfrac12\cdot96=-48$. The left-hand side comes
from \eqref{eq:moonshinecomputed:chi2Y}: $F_2^{(2,1)}(1)=1\cdot(1\cdot[2>1]-2\cdot[1>2])=1$,
so $[q^1]\chi_2^{A_1}=1\cdot(-24)-24\cdot1=-48$. The two agree, and the constant terms agree
as well ($\mathrm{rk}=1$ against $-\tfrac12\cdot(-2)=1$). The general statement, for the six
lambencies with $X$ a power of a single $A$ component, is checked in Lean through $q^6$.
\end{example}

HMN relate $\chi_2^Y$ to umbral moonshine for the Niemeier root system
$X=Y^{24/\mathrm{rk}(Y)}$ \source{HMN eqs.~(4.5), (4.9), ll.~1229--1266}:
\begin{equation}\label{eq:moonshinecomputed:umbral}
  \chi_2^Y=-\frac{\mathrm{rk}(Y)}2\,\chi_2^X,\qquad \chi_2^X=\sum_rH^X_r\,S_{m,r}.
\end{equation}
For $Y=A_{\ell-1}$ the umbral forms can be built from theta functions exactly as in
\cref{sec:moonshinecomputed:shadow}, now with CDH's $Z^{(\ell)}$ for
$\ell=2,3,4,5,7,13$ \source{CDH eqs.~(2.49)--(2.51), (2.62), ll.~1258--1285, 1397}; the
quantity $\chi_2^X$ is the first Taylor coefficient at $z=0$ of the finite part. Lean checks
\eqref{eq:moonshinecomputed:umbral} for all six values of $\ell$ through $q^6$, and that
pairing $\ell=3$ with $Y=A_3$ instead of $A_2$ fails.

HMN observe that for the ``umbral'' $Y$ (those with $\mathrm{rk}(Y)\mid24$) every coefficient
of $\chi_2^Y$ is divisible by $\mathrm{rk}(Y)$, and write that ``it is not clear why this
should be true for the coefficients of higher powers of $q$'' \source{HMN §4, ll.~1280--1296}.
Within their own formula \eqref{eq:moonshinecomputed:chi2Y} the reason is elementary. For
$n\ge1$,
\begin{equation}
  [q^n]\chi_2^Y=-24\bigl(\mathrm{rk}(Y)\,\sigma_1(n)+F_2^Y(n)\bigr),\qquad F_2^Y(n)\in\Z,
\end{equation}
so $\mathrm{rk}(Y)\mid24$ forces divisibility of every coefficient. Conversely, for
$A_{k-1}$ one has $F_2^Y(1)=1$ and $[q^1]\chi_2^Y=-24(\mathrm{rk}+1)$, divisible by
$\mathrm{rk}$ only if $\mathrm{rk}\mid24$; for $D_{j+1}$, $F_2^Y(1)=-1$ gives
$-24(\mathrm{rk}-1)$ with the same conclusion. Lean proves both directions for every
$A_{k-1}$ and every $D_{j+1}$, and checks $E_6,E_8$ (divisible) and $E_7$ (not). The
divisible $Y$ are exactly HMN's umbral list $A_1,A_2,A_3,A_4,A_6,A_8,A_{12},A_{24},D_4,D_6,
D_8,D_{12},D_{24},E_6,E_8$. Why the factor in \eqref{eq:moonshinecomputed:chi2Y} is $24$ is
not addressed here; reading it as $\chi(\mathrm{K3})$ is \TC.

\begin{pitfallbox}[A sign in the source]
HMN's eq.~(4.2) \source{HMN eq.~(4.2), l.~1194} prints $\chi_2^Y=-\mathrm{rk}(Y)E_2+24F_2^Y$,
the opposite overall sign of their eq.~(2.36). The printed series (4.10)--(4.11) start with
$+\mathrm{rk}(Y)$, as does the count of $\mathrm{rk}(Y)$ massless vector multiplets
\source{HMN §4, ll.~1257--1260}; they fix the sign of (2.36), which is the one used here.
\end{pitfallbox}

\subsection{What is not available: a reflection bridge}

The vanishing $(-2)$-cycles of $Y$ suggest a link with the reflections in $(-2)$-vectors
that \cref{ch:mukai} and Stream~2 formalize. HMN's lattice structure, however, enters
through the Niemeier lattice\index{Niemeier lattice} $L_X$ and its Weyl group,\index{Weyl group}
$G^X=\mathrm{Aut}(L_X)/W_X$ \source{HMN eq.~(1.10), l.~277}, a rank-$24$ even unimodular
lattice that this project does not formalize. Nothing in the computations of this section
involves a reflection or a lattice isometry; $\mathrm{rk}(Y)$ and the Coxeter number enter
only as numbers. The bridge is therefore recorded as \emph{not available}, with that
reason, rather than claimed.

\section{What the machine has checked}\label{sec:moonshinecomputed:lean}

The files are in the library \lean{DualScaleMoonshine}. All theorems quoted depend only on
\lean{propext}, \lean{Classical.choice} and \lean{Quot.sound} (most on \lean{propext}
alone); the heavy ones are closed by \lean{decide} or \lean{decide +kernel}, which have the
kernel evaluate the decision procedure.

\begin{leanbox}[In Lean: $H$ computed equals the table]
\begin{leancode}
def numer (N : ℕ) : List ℤ :=
  (List.range (N + 1)).map fun n => if n = 0 then -2 else 48 * (sigma n + f2Coeff n)

def hComputed (N : ℕ) : List ℤ := divTrunc N (numer N) (eta3 N)

theorem eta3_jacobi : eta3 9 = [1, -3, 0, 5, 0, 0, -7, 0, 0, 0] := by ...

theorem hComputed_eq_table : hComputed 9 = hFromTable := by ...
\end{leancode}
(\lean{DualScaleMoonshine/QSeries.lean}.) \textbf{Reading.} \lean{hFromTable} is
$[-2,2A_1,\dots,2A_9]$ from EOT's table; the list computed by the recursion
\eqref{eq:moonshinecomputed:div} from $-2E_2+48F_2$ and $\prod(1-q^n)^3$ equals it through
$q^9$. \textbf{Proof idea.} Kernel evaluation of two lists of ten integers.
\textbf{Tier} \TA\ for the lists; \TL\ that the formula is the K3 mock modular form.
\end{leanbox}

\begin{leanbox}[In Lean: twined series, with newforms and controls]
\begin{leancode}
theorem twined_2A : twined24 9 8 2 (-16) = table2A.map (24 * ·) := by ...

theorem twined_11A : twinedD dataF11A = scaled dataF11A table11A := by ...
theorem twined_23AB : twinedD dataF23AB = scaled dataF23AB table23AB := by ...

theorem twined_11A_needs_newform :
    twinedD (5, 2, lambda24 9 11 (-2)) ≠ scaled dataF11A table11A := by ...
\end{leancode}
(\lean{DualScaleMoonshine/Twining.lean}, \lean{TwiningAll.lean}.) \textbf{Reading.}
$24H_{2A}$ computed from $\chi_{2A}=8$ and $F_{2A}=-16\Lambda_2$ equals $24$ times CDH's
printed $2A$ column through $q^9$; likewise $24D\,H_g$ for $11A$ ($D=5$) and $23AB$
($D=11$), whose $F_g$ contain newforms; dropping $f_{11}$ breaks $11A$. All $21$ columns have
such a theorem. \textbf{Tier} \TA\ for the arithmetic, \TL\ for the formulas of CDH.
\end{leanbox}

\begin{leanbox}[In Lean: the character table and the traces]
\begin{leancode}
theorem colNorm_eq_cent : (List.range 26).map colNorm = cent.map (·, 0) := by ...
theorem class_equation : classSize.sum = 244823040 := by ...
theorem gram_ok : gramOK charTab = true := by ...

theorem trace_eq_twined_coeff_all :
    ∀ j : Fin 26, eotMult.map (traceQ j) =
      (((computedSeries.getD j []).drop 1).take 7).map (·, 0) := by ...
\end{leancode}
(\lean{DualScaleMoonshine/CharactersAll.lean}.) \textbf{Reading.} The column norms of the
table (entries in $\Z[b_n]$ stored as pairs) are the listed centralizer orders; the class
sizes add to $|\Mtf|$; all $26\times26$ orthogonality relations hold; and at every class
$j$, the traces on EOT's representations at levels $1$--$7$ (doubled, as in
\lean{eotMult}) have zero irrational part and equal the computed twined coefficients.
\textbf{Tier} \TA; that the table is the character table of $\Mtf$ is \TL.
\end{leanbox}

\begin{leanbox}[In Lean: the modules $K_0,\dots,K_9$ decomposed]
\begin{leancode}
theorem moonshine_modules_decompose :
    (List.range 10).all (fun n => (List.range 26).all fun i =>
      let v := moonshineInner n i
      v.2.1 == 0 && v.2.2.1 == 0 && v.2.2.2 == 0 &&
        v.1 == 244823040 * ((table48.getD n []).getD i 0)) = true := by ...
\end{leancode}
(\lean{DualScaleMoonshine/Decompositions.lean}.) \textbf{Reading.} For $n\le9$ and every
irreducible $i$, the inner product \eqref{eq:moonshinecomputed:mult} computed from the
twined series has vanishing irrational parts and equals $|\Mtf|$ times CDH's Table~48 entry,
so every multiplicity is an integer equal to the printed one. A companion theorem,
\lean{eot\_level7\_proposal\_inconsistent}, shows the traces of the older level-$7$
proposal differ from the computed series. \textbf{Tier} \TA.
\end{leanbox}

\begin{leanbox}[In Lean: the forger's test]
\begin{leancode}
theorem lock_at_identity :
    coeff untwined 1 = 90 ∧ coeff untwined 2 = 462 ∧ literalLock untwined ∧
      coeff untwined 2 * 60 = 27720 := by ...

theorem ratio_fails_at_every_class :
    ∀ j : Fin 26, j.val ≠ 0 →
      ¬ ratioTwines ((computedSeries.getD j []).map (24 * ·)) := by ...
\end{leancode}
(\lean{DualScaleMoonshine/ForgerTest.lean}, \lean{CharactersAll.lean}.) \textbf{Reading.}
The lock holds for the untwined series and its ratio form fails for the twined series of
every one of the $25$ non-identity classes. \textbf{Tier} \TA\ for both; the conclusion
``numerology'' is the reading the twining criterion gives (\TC\ as a methodological
judgement).
\end{leanbox}

\begin{leanbox}[In Lean: the 24 of the shadow]
\begin{leancode}
theorem ellipticGenus_z0 :
    (ellipticGenus 9).map LP.eval1 = [24, 0, 0, 0, 0, 0, 0, 0, 0, 0] := by ...

theorem decomposition : Ser.eqB (starLHS 9) (starRHS 9 24 (hComputed 9)) = true := by ...

theorem decomposition_pins_chi :
    Ser.eqB (starLHS 9) (starRHS 9 23 (hComputed 9)) = false ∧
      Ser.eqB (starLHS 9) (starRHS 9 25 (hComputed 9)) = false := by ...
\end{leancode}
(\lean{DualScaleMoonshine/Shadow.lean}.) \textbf{Reading.} The elliptic genus computed from
theta products is $24$ at $y=1$ at every order through $q^9$; the decomposition
\eqref{eq:moonshinecomputed:polar}, cleared of denominators, holds with polar multiplicity
$24$ and finite part the computed $H$, and fails for $23$ and $25$. \textbf{Tier} \TA; the
word ``shadow'' is \TL.
\end{leanbox}

\begin{leanbox}[In Lean: the little-string index]
\begin{leancode}
def chi2Y (Ys : List (ℕ × ℕ)) (rk : ℤ) (n : ℕ) : ℤ :=
  if n = 0 then rk else rk * (-24 * sigma n) - 24 * (Ys.map fun p => f2kd p.1 p.2 n).sum

theorem hmn_A3_printed : chi2YSer (yA 4) 3 5 = [3, -96, -288, -384, -576, -360] := by ...
theorem hmn_k2_eq_moonshine : (chi2YSer (yA 2) 1 9).map (-2 * ·) = numer 9 := by ...
theorem hmn_umbral_relation_3 :
    (chi2YSer (yA 3) 2 6).map (4 * ·) = (thetaPsiF 6 3).map (2 * ·) := by ...
\end{leancode}
(\lean{DualScaleMoonshine/HMNBridge.lean}.) \textbf{Reading.} The first theorem reproduces
HMN's printed $A_3$ series; the second says the $k=2$ index times $-2$ is the moonshine
numerator; the third is \eqref{eq:moonshinecomputed:umbral} for $\ell=3$ through $q^6$, with
the umbral side built from theta functions (the other lambencies have their own theorems).
The divisibility statements \lean{divisible\_of\_dvd\_24}, \lean{A\_divisible\_iff} and
\lean{D\_divisible\_iff} are ordinary proofs quantified over all $n$ and all $k$ (resp.\
$j$); their statements use the divisibility sign of Lean, which this book's fonts do not
typeset, so we give them in words: if $\mathrm{rk}\mid24$ then $\mathrm{rk}$ divides every
coefficient; for $A_{k-1}$ ($k\ge2$) and $D_{j+1}$ ($j\ge3$), all coefficients are divisible
by the rank if and only if the rank divides $24$. \textbf{Tier} \TA; the physical meaning
of $\chi_2^Y$ is \TL.
\end{leanbox}

\begin{remark}[What the checks cost]
Exactness is not free. The heaviest file, the little-string bridge, takes about a quarter of
an hour of kernel time: the products of eta quotients and the sixteen-fold Newton
recursions are evaluated on integers, with no floating point anywhere. Three theorems in it
use \lean{decide +kernel}, which skips the elaborator's own evaluation and lets the kernel
do the arithmetic directly; this adds no axiom and is not the compiled evaluation that the
project's rules forbid. A reader who wants to reproduce a single column should compute it
with the recursion \eqref{eq:moonshinecomputed:div} in any language: the numbers are small.
\end{remark}

\begin{center}\small
\begin{tabular}{p{7.3cm}ll}
\hline
Claim & Tier & Where\\ \hline
$H$ from $(-2E_2+48F_2)/\eta^3$ equals $2A_n$, $n\le9$ & \TA & \lean{hComputed\_eq\_table}\\
$21$ twined columns equal CDH Table~20 through $q^9$ & \TA & \lean{twined\_2A}, \dots\\
traces $=$ twined coefficients, all classes, $n\le7$ & \TA & \lean{trace\_eq\_twined\_coeff\_all}\\[2pt]
$K_0,\dots,K_9$ decompose as CDH Table~48 & \TA & \lean{moonshine\_modules\_decompose}\\
$27720$ lock fails twining at $25$ classes & \TA & \lean{ratio\_fails\_at\_every\_class}\\
$Z(\tau,0)=24=$ polar multiplicity; $23,25$ fail & \TA & \lean{decomposition\_pins\_chi}\\
little-string index $=$ umbral form, $\ell\le13$ & \TA & \lean{hmn\_umbral\_relation\_3}\\
divisibility iff $\mathrm{rk}\mid24$ ($A$, $D$) & \TA & \lean{A\_divisible\_iff}\\
$H$ is mock modular with shadow $24\eta^3$ & \TL & CDH, HMN\\
$\Mtf$ acts on a module with these traces & \TL & Gannon (via CDH)\\
$24$ in \eqref{eq:moonshinecomputed:chi2Y} is $\chi(\mathrm{K3})$ & \TC & ---\\
reflection bridge to Stream~2 & not available & \cref{sec:moonshinecomputed:hmn}\\ \hline
\end{tabular}
\end{center}

\begin{summarybox}
\begin{itemize}[leftmargin=*]
\item Truncated $q$-series are finite lists of integers; products and quotients by series
  with constant term $1$ are exact, so identities through $q^N$ are decidable.
\item $q^{1/8}H=(-2E_2+48F_2)/\prod(1-q^n)^3=-2+90q+462q^2+1540q^3+\dots$ reproduces
  $2A_1,\dots,2A_9$; the twined $H_g=\frac{\chi_g}{24}H+F_g/\eta^3$ reproduces all $21$ columns
  of CDH's table, including those needing newforms.
\item The character table over $\Z[b_7],\Z[b_{15}],\Z[b_{23}]$ is guarded by orthogonality
  and the class equation (up to relabelling); traces equal computed coefficients at every
  class, and $K_0,\dots,K_9$ decompose with the printed multiplicities.
\item Forger's test: the $27720$ lock fails at all $25$ non-identity classes; the $24$
  twines into the character of $\mathbf1\oplus\mathbf{23}$.
\item $Z(\tau,0)=24$ and the polar multiplicity of the elliptic genus are the same integer,
  and only $24$ works; the shadow statement itself is \TL.
\item HMN's little-string index reproduces Mathieu moonshine at $k=2$ and umbral moonshine
  for $\ell=3,4,5,7,13$; its coefficients are divisible by $\mathrm{rk}(Y)$ iff
  $\mathrm{rk}(Y)\mid24$, for elementary reasons; the reflection bridge is not available.
\end{itemize}
\end{summarybox}

\section*{Exercises}
\addcontentsline{toc}{section}{Exercises}

\begin{exercise}[$\star$]\label{exo:moonshinecomputed:h4}
Continue the recursion of \cref{sec:moonshinecomputed:worked} to $q^4$. You need
$[q^4]F_2=1$, $\sigma_1(4)=7$ and $p_4=0$. Check that you obtain $2A_4=4554$.
\end{exercise}

\begin{exercise}[$\star$]\label{exo:moonshinecomputed:2Aq2}
For the class $2A$, compute $[q^2]$ of $24H_{2A}$ from $8s+24F_{2A}$ and the recursion,
and verify that $[q^{15/8}]H_{2A}=14=2\chi_{\mathbf{231}}(2A)$.
\end{exercise}

\begin{exercise}[$\star$]\label{exo:moonshinecomputed:lock}
Show that the literal lock $A_2(g)\cdot60=4A_1(g)\cdot77$ fails at $3A$, using
$A_1(3A)=0$, $A_2(3A)=-3$ (half the entries $0,-6$ of CDH's $3A$ column). Why is the ratio
form of the lock the fairer test?
\end{exercise}

\begin{exercise}[$\star\star$]\label{exo:moonshinecomputed:b7}
Verify $b_7^2=-b_7-2$ and $b_7\overline{b_7}=2$ from the definition, and show that the
rule $(a,b)\cdot(c,d)=(ac-kbd,\ ad+bc-bd)$ with $k=(n+1)/4$ is multiplication in
$\Z[b_n]$. Why is $\Z[b_n]$ closed under complex conjugation?
\end{exercise}

\begin{exercise}[$\star\star$]\label{exo:moonshinecomputed:A7}
Show that $10395+2\cdot5796+5544+3312=30843=A_7$, so the older proposal has the right
dimension. Using the $2A$ column of the character table, compute the trace of $2A$ on the
doubled module and compare with $[q^{55/8}]H_{2A}=-138$.
\end{exercise}

\begin{exercise}[$\star\star$]\label{exo:moonshinecomputed:D4}
For $Y=D_4$ ($m=6$, $j=3$) compute $F_2^{(6,1)}(1)$ and $F_2^{(6,3)}(1)$ from
\eqref{eq:moonshinecomputed:chi2Y}, and verify $[q^1]\chi_2^{D_4}=-72$. Is it divisible by
$\mathrm{rk}(D_4)=4$? Repeat for $D_5$ and explain the difference.
\end{exercise}

\begin{exercise}[$\star\star$]\label{exo:moonshinecomputed:lean1}
In Lean, state and prove by \lean{decide} that the level-$9$ multiplicities of
\lean{table48} sum, weighted by the dimensions in \lean{chi1A}, to $2A_9=265650$. Which
theorem of this chapter makes your statement redundant, and why is it still a useful
check on the transcription of Table~48?
\end{exercise}

\begin{exercise}[$\star\star\star$]\label{exo:moonshinecomputed:lean2}
Write a Lean definition of the ratio relation for an arbitrary pair of levels $(a,b)$ and
prove, by \lean{decide} on \lean{computedSeries}, that $A_3(g)/A_1(g)=A_3/A_1$ also fails at
some class. Then formulate a relation between the $A_n$ that \emph{does} survive twining at
every class (hint: a decomposition of $K_6$) and prove it.
\end{exercise}

\begin{exercise}[$\star\star$]\label{exo:moonshinecomputed:etaq}
Check that the eta quotient $-2\eta(\tau)^8/\eta(2\tau)^4$ of the class $2B$
\source{CDH Table~3, ll.~2882--2915} satisfies $24\mid\sum_kke_k$, and compute its first two
coefficients from \eqref{eq:moonshinecomputed:etaq}. Compare with $24\Lambda_2-8\Lambda_4$,
the other form in which CDH print the same $F_{2B}$.
\end{exercise}

\begin{exercise}[$\star\star\star$]\label{exo:moonshinecomputed:shadow}
Clear the denominators in \eqref{eq:moonshinecomputed:polar} for $m=2$: with
$E=\prod(1-q^n)$, $B_1=\prod(1-yq^n)(1-y^{-1}q^n)$ and $B_2=\prod(1-y^2q^n)(1-y^{-2}q^n)$,
show that $\Psi_{1,1}=\frac{y+1}{y-1}\,E^2B_2/B_1^2$, and check the first two coefficients
$(y+1)/(y-1)-(y^2-y^{-2})q$ against \source{CDH eq.~(2.20), ll.~670--675}.
\end{exercise}

\section*{Further reading}
\addcontentsline{toc}{section}{Further reading}
\begin{itemize}[leftmargin=*]
\item The closed formula for $H$: \source{CH eqs.~(3.5)--(3.6), ll.~797--815}; the
  twined series, the forms $F_g$, the twisted Euler characters and the printed columns:
  \source{CDH eq.~(4.18), ll.~2931--2936}, \source{CDH Table~3, ll.~2882--2915},
  \source{CDH Table~14, ll.~4349--4364}, \source{CDH Table~20, ll.~4494--4512}; the
  newforms: \source{CDH App.~A, ll.~4100--4137}.
\item The character table and the decompositions: \source{CDH Table~8, ll.~4170--4209},
  \source{CDH Table~48, ll.~5305--5330}; the level-$7$ discussion: \source{Cheng,
  ll.~533--549}.
\item Meromorphic Jacobi forms, polar and finite parts, and the shadow:
  \source{CDH §2.3, ll.~657--768}; the umbral Jacobi forms: \source{CDH §2.5,
  ll.~1258--1397}; the shadow $24\eta^3$ and the mixed mock form $\eta^3H$:
  \source{HMN §1, ll.~244--270}.
\item The little-string index, its ADE combinations and the umbral relation:
  \source{HMN §2.1--2.2, ll.~587--736}, \source{HMN §4, ll.~1179--1300}.
\item The observation, twining genera and the older decompositions are in
  \cref{ch:moonshine}; the group $\Mtf$ in \cref{ch:m24}; the elliptic genus and its
  characters in \cref{ch:ellgenus}; the dyon partition function built from the same
  elliptic genus in \cref{ch:dyons}. The project record is
  \texttt{docs/STREAM4\_WORKFLOW.md}.
\end{itemize}
